\chapter{Conclusion}
\label{chap:conclusion}

The proposed distributed node architecture was validated, and separating the audio controller from the RF synchronization controller proved to be an appropriate design decision. The audio subsystem fulfilled its core requirements: it achieved \SI{48}{\kilo\hertz} acquisition, full-duplex playback and recording, local WAV storage, and exponential-sine-sweep impulse-response recovery. Crucially, its round-trip latency was deterministic and correctable; the system does not eliminate latency, but makes it fixed and repeatable (about \SI{39.7}{\milli\second} electrical and \SI{42.9}{\milli\second} acoustic, stable to within one sample), which is the key requirement for reconstructing the impulse response in post-processing. The analog front end fulfilled the measurement-interface requirements, with the IEPE constant-current source, gain-ranged signal paths, signal conditioning, and headroom all validated for measurement microphones. On the wireless side, the synchronization requirement was met at the microsecond scale: the event-based master--slave protocol achieved repeatable trigger alignment without continuous clock synchronization, and, with external antennas, the RF link demonstrated long-range feasibility at distances on the order of \SI{1}{\kilo\meter}, although the achievable range depends strongly on antenna installation and environment.
% TODO: confirm the "on the order of 1 km" figure against the measured RF range results, and state the actual achieved distance.

Integrated into a working node, the system produced physically consistent reverberation-time measurements---agreeing with a reference analysis routine---and two-room sound-insulation results that met the ISO~16283 margin at low frequencies, confirming that the distributed, wirelessly synchronized approach is a valid basis for standard room-acoustic metrology. The power architecture generated the required rails, but integration noise limited the autonomous node: the main issue was not the audio chain itself---whose front-end thermal noise was deliberately kept low---but the custom supply and grounding path, particularly toward the audio-controller board. This supply-conducted disturbance, dominated by a \SI{50}{\hertz} mains-hum component coupled as radiated interference through the unshielded interconnect, set the usable dynamic range at about \SI{66}{\decibel}---below what the codec and front end could otherwise provide---while powering the node from a battery bank alone recovered \SIrange{21}{25}{\decibel} of that margin and removed the hum. The battery sizing was reasonable and the measured consumption supports the energy budget, but the autonomy was only partially validated, as a complete discharge profile was not characterized.

Overall, the system works as an experimental acoustic measurement node. It is not yet a fully optimized field instrument, but it validates the main engineering concept: a deterministic, wirelessly synchronized, distributed architecture for impulse-response measurement. The remaining challenges are primarily integration challenges. The next iteration should therefore focus on reducing power-supply noise, improving grounding and PCB layout, characterizing the antennas and long-range timing under field conditions, and adding low-power operating modes, together with a second analog-front-end board to exploit the shared time base for simultaneous dual-room capture.
